% --- Archon physics formalization source begin ---
% archon:physics
% archon:covers PhyXMiniProblems/problem_phyx_mini_0966.lean
% archon:source-report reports/phyx_mini/problem_phyx_mini_0966.source.json

\chapter{Physics problem phyx\_mini\_0966}
\label{ch:PhyXMiniProblems_problem_phyx_mini_0966}

\paragraph{Problem source.}
\#\# Physical scenario

Segment \textbackslash{}( DA \textbackslash{}) is an arc of a circle with radius \textbackslash{}( 20.0\textbackslash{},\textbackslash{}text\{cm\} \textbackslash{}), and point \textbackslash{}( P \textbackslash{}) is at its center of curvature. Segments \textbackslash{}( CD \textbackslash{}) and \textbackslash{}( AB \textbackslash{}) are straight lines of length \textbackslash{}( 10.0\textbackslash{},\textbackslash{}text\{cm\} \textbackslash{}) each.

\#\# Question

Calculate the magnetic field (magnitude) at a point \textbackslash{}( P \textbackslash{}) due to a current \textbackslash{}( I = 12.0\textbackslash{},\textbackslash{}text\{A\} \textbackslash{}) in the wire. Segment \textbackslash{}( BC \textbackslash{}) is an arc of a circle with radius \textbackslash{}( 30.0\textbackslash{},\textbackslash{}text\{cm\} \textbackslash{}), and point \textbackslash{}( P \textbackslash{}) is at the center of curvature of the arc.

\#\# Answer choices

- A: A: 4.29μT

- B: B: 4.19μT

- C: C: 6.28μT

- D: D: 2.09μT

\#\# Figure caption (auxiliary; use the image as primary evidence)

The image depicts a geometric diagram involving a segment of a circle. Here are the components and their relationships:

- The figure consists of two concentric arcs centered at a point \textbackslash{}( P \textbackslash{}).
- The arcs form part of a circle of a sector with a central angle of \textbackslash{}( 120\textasciicircum{}\textbackslash{}circ \textbackslash{}).
- The sector is bounded by two radii, \textbackslash{}( PA \textbackslash{}) and \textbackslash{}( PB \textbackslash{}), and two curves, \textbackslash{}( CD \textbackslash{}) and \textbackslash{}( AB \textbackslash{}).
- Points \textbackslash{}( C \textbackslash{}) and \textbackslash{}( D \textbackslash{}) are on the inner arc, while \textbackslash{}( A \textbackslash{}) and \textbackslash{}( B \textbackslash{}) are on the outer arc with a clockwise orientation.
- There are purple arrows labeled with \textbackslash{}( I \textbackslash{}) along the segments \textbackslash{}( CD \textbackslash{}), \textbackslash{}( AB \textbackslash{}), and the arc between them, indicating the direction of some vector or current.
- The angle \textbackslash{}( 120\textasciicircum{}\textbackslash{}circ \textbackslash{}) is marked between the lines \textbackslash{}( PA \textbackslash{}) and \textbackslash{}( PB \textbackslash{}).
- The arcs and radial lines form the sector or slice of a circle with a gap between \textbackslash{}( CD \textbackslash{}) and \textbackslash{}( AB \textbackslash{}).

\#\# Dataset metadata

Magnetism; Physical Model Grounding Reasoning, Spatial Relation Reasoning

\paragraph{Recorded answer/context.}
B: B: 4.19μT

\paragraph{Figure/image path.}
/root/proposal\_for\_physic/science-mango/phyx\_mini\_run/phyx\_data/test\_image/966.png

\paragraph{Formalization target.}
create a compiling Lean file with sorry bodies at `PhyXMiniProblems/problem\_phyx\_mini\_0966.lean`. The Lean declarations must preserve the physical quantities, dimensions or dimensional roles, figure labels, governing-law hypotheses, and final relation expressed by this problem.
Use Mathlib/Physlib names found through LeanExplore where available. If a physics API is missing, introduce faithful local abstractions rather than scalar placeholder aliases.

\begin{theorem}[Physics formalization target]
\label{thm:physics:phyx_mini_0966:target}
\lean{PhyXMiniProblems.ProblemPhyXMini0966.problem_phyx_mini_0966}
\uses{def:physics:phyx-mini-0966:phyxminiproblems-problemphyxmini0966-magneticfluxdensityinmicroteslas, def:physics:phyx-mini-0966:phyxminiproblems-problemphyxmini0966-annularsectorwiresetup, def:physics:phyx-mini-0966:phyxminiproblems-problemphyxmini0966-matchesannularsectorproblemdata, def:physics:phyx-mini-0966:phyxminiproblems-problemphyxmini0966-matchesprimaryannularsectorfigure, def:physics:phyx-mini-0966:phyxminiproblems-problemphyxmini0966-hasannularsectorgeometry, def:physics:phyx-mini-0966:phyxminiproblems-problemphyxmini0966-hasphysicalannularsectorparameters, def:physics:phyx-mini-0966:phyxminiproblems-problemphyxmini0966-usesstandardvacuumpermeability, def:physics:phyx-mini-0966:phyxminiproblems-problemphyxmini0966-satisfiessteadycurrentfieldmodel, def:physics:phyx-mini-0966:phyxminiproblems-problemphyxmini0966-satisfiescirculararcbiotsavartlaw, def:physics:phyx-mini-0966:phyxminiproblems-problemphyxmini0966-satisfiesradialconnectorbiotsavartlaw, def:physics:phyx-mini-0966:phyxminiproblems-problemphyxmini0966-satisfiesmagneticfieldsuperposition, def:physics:phyx-mini-0966:phyxminiproblems-problemphyxmini0966-calibratesmagneticfluxdensitymagnitudeatp, def:physics:phyx-mini-0966:phyxminiproblems-problemphyxmini0966-displayedmagneticfieldinmicroteslas, def:physics:phyx-mini-0966:phyxminiproblems-problemphyxmini0966-roundstonearesthundredthmicrotesla, def:physics:phyx-mini-0966:phyxminiproblems-problemphyxmini0966-isuniqueclosestdisplayedanswer}
The assigned autoformalize agent should translate this physics problem into Lean declarations in the covered file, with theorem and lemma proof bodies written as `by sorry`.
\end{theorem}
\begin{proof}
This is an autoformalization task, not a proof task. Produce statements that can later be proved by the physics prover without weakening the physical model.
\end{proof}

% --- Archon declaration topology begin ---
\section{Declaration topology}
\begin{definition}[electric Current Dimension]
  \label{def:physics:phyx-mini-0966:phyxminiproblems-problemphyxmini0966-electriccurrentdimension}
  \lean{PhyXMiniProblems.ProblemPhyXMini0966.electricCurrentDimension}
  Electric current has dimension charge per unit time
\end{definition}
\begin{proof}
  This is the declared model definition of electric Current Dimension.
\end{proof}
\begin{definition}[magnetic Permeability Dimension]
  \label{def:physics:phyx-mini-0966:phyxminiproblems-problemphyxmini0966-magneticpermeabilitydimension}
  \lean{PhyXMiniProblems.ProblemPhyXMini0966.magneticPermeabilityDimension}
  Magnetic permeability has dimension M L C⁻²
\end{definition}
\begin{proof}
  This is the declared model definition of magnetic Permeability Dimension.
\end{proof}
\begin{definition}[magnetic Flux Density Dimension]
  \label{def:physics:phyx-mini-0966:phyxminiproblems-problemphyxmini0966-magneticfluxdensitydimension}
  \lean{PhyXMiniProblems.ProblemPhyXMini0966.magneticFluxDensityDimension}
  Magnetic flux density (tesla) has dimension M T⁻¹ C⁻¹
\end{definition}
\begin{proof}
  This is the declared model definition of magnetic Flux Density Dimension.
\end{proof}
\begin{definition}[Length Quantity]
  \label{def:physics:phyx-mini-0966:phyxminiproblems-problemphyxmini0966-lengthquantity}
  \lean{PhyXMiniProblems.ProblemPhyXMini0966.LengthQuantity}
  A nonnegative, unit-independent physical length
\end{definition}
\begin{proof}
  This is the declared model definition of Length Quantity.
\end{proof}
\begin{definition}[Electric Current Quantity]
  \label{def:physics:phyx-mini-0966:phyxminiproblems-problemphyxmini0966-electriccurrentquantity}
  \lean{PhyXMiniProblems.ProblemPhyXMini0966.ElectricCurrentQuantity}
  \uses{def:physics:phyx-mini-0966:phyxminiproblems-problemphyxmini0966-electriccurrentdimension}
  A nonnegative, unit-independent steady-current magnitude
\end{definition}
\begin{proof}
  This is the declared model definition of Electric Current Quantity.
\end{proof}
\begin{definition}[Magnetic Permeability Quantity]
  \label{def:physics:phyx-mini-0966:phyxminiproblems-problemphyxmini0966-magneticpermeabilityquantity}
  \lean{PhyXMiniProblems.ProblemPhyXMini0966.MagneticPermeabilityQuantity}
  \uses{def:physics:phyx-mini-0966:phyxminiproblems-problemphyxmini0966-magneticpermeabilitydimension}
  A nonnegative, unit-independent magnetic permeability
\end{definition}
\begin{proof}
  This is the declared model definition of Magnetic Permeability Quantity.
\end{proof}
\begin{definition}[Magnetic Flux Density Quantity]
  \label{def:physics:phyx-mini-0966:phyxminiproblems-problemphyxmini0966-magneticfluxdensityquantity}
  \lean{PhyXMiniProblems.ProblemPhyXMini0966.MagneticFluxDensityQuantity}
  \uses{def:physics:phyx-mini-0966:phyxminiproblems-problemphyxmini0966-magneticfluxdensitydimension}
  A nonnegative, unit-independent magnetic-flux-density magnitude
\end{definition}
\begin{proof}
  This is the declared model definition of Magnetic Flux Density Quantity.
\end{proof}
\begin{definition}[length In Meters]
  \label{def:physics:phyx-mini-0966:phyxminiproblems-problemphyxmini0966-lengthinmeters}
  \lean{PhyXMiniProblems.ProblemPhyXMini0966.lengthInMeters}
  \uses{def:physics:phyx-mini-0966:phyxminiproblems-problemphyxmini0966-lengthquantity}
  Coherent-SI metre readout of a physical length
\end{definition}
\begin{proof}
  This is the declared model definition of length In Meters.
\end{proof}
\begin{definition}[length In Centimeters]
  \label{def:physics:phyx-mini-0966:phyxminiproblems-problemphyxmini0966-lengthincentimeters}
  \lean{PhyXMiniProblems.ProblemPhyXMini0966.lengthInCentimeters}
  \uses{def:physics:phyx-mini-0966:phyxminiproblems-problemphyxmini0966-lengthquantity, def:physics:phyx-mini-0966:phyxminiproblems-problemphyxmini0966-lengthinmeters}
  Centimetre readout used by the three printed length values
\end{definition}
\begin{proof}
  This is the declared model definition of length In Centimeters.
\end{proof}
\begin{definition}[current In Amperes]
  \label{def:physics:phyx-mini-0966:phyxminiproblems-problemphyxmini0966-currentinamperes}
  \lean{PhyXMiniProblems.ProblemPhyXMini0966.currentInAmperes}
  \uses{def:physics:phyx-mini-0966:phyxminiproblems-problemphyxmini0966-electriccurrentquantity}
  Coherent-SI ampere readout of the current magnitude
\end{definition}
\begin{proof}
  This is the declared model definition of current In Amperes.
\end{proof}
\begin{definition}[permeability In Tesla Meters Per Ampere]
  \label{def:physics:phyx-mini-0966:phyxminiproblems-problemphyxmini0966-permeabilityinteslametersperampere}
  \lean{PhyXMiniProblems.ProblemPhyXMini0966.permeabilityInTeslaMetersPerAmpere}
  \uses{def:physics:phyx-mini-0966:phyxminiproblems-problemphyxmini0966-magneticpermeabilityquantity}
  Coherent-SI permeability readout in tesla-metres per ampere
\end{definition}
\begin{proof}
  This is the declared model definition of permeability In Tesla Meters Per Ampere.
\end{proof}
\begin{definition}[magnetic Flux Density In Teslas]
  \label{def:physics:phyx-mini-0966:phyxminiproblems-problemphyxmini0966-magneticfluxdensityinteslas}
  \lean{PhyXMiniProblems.ProblemPhyXMini0966.magneticFluxDensityInTeslas}
  \uses{def:physics:phyx-mini-0966:phyxminiproblems-problemphyxmini0966-magneticfluxdensityquantity}
  Coherent-SI tesla readout of a magnetic-flux-density magnitude
\end{definition}
\begin{proof}
  This is the declared model definition of magnetic Flux Density In Teslas.
\end{proof}
\begin{definition}[magnetic Flux Density In Microteslas]
  \label{def:physics:phyx-mini-0966:phyxminiproblems-problemphyxmini0966-magneticfluxdensityinmicroteslas}
  \lean{PhyXMiniProblems.ProblemPhyXMini0966.magneticFluxDensityInMicroteslas}
  \uses{def:physics:phyx-mini-0966:phyxminiproblems-problemphyxmini0966-magneticfluxdensityquantity, def:physics:phyx-mini-0966:phyxminiproblems-problemphyxmini0966-magneticfluxdensityinteslas}
  Microtesla readout used by the multiple-choice answers
\end{definition}
\begin{proof}
  This is the declared model definition of magnetic Flux Density In Microteslas.
\end{proof}
\begin{definition}[Spatial Vector]
  \label{def:physics:phyx-mini-0966:phyxminiproblems-problemphyxmini0966-spatialvector}
  \lean{PhyXMiniProblems.ProblemPhyXMini0966.SpatialVector}
  A coherent-SI three-dimensional spatial vector
\end{definition}
\begin{proof}
  This is the declared model definition of Spatial Vector.
\end{proof}
\begin{definition}[Figure Point]
  \label{def:physics:phyx-mini-0966:phyxminiproblems-problemphyxmini0966-figurepoint}
  \lean{PhyXMiniProblems.ProblemPhyXMini0966.FigurePoint}
  The five point labels printed in the primary image
\end{definition}
\begin{proof}
  This is the declared model definition of Figure Point.
\end{proof}
\begin{definition}[Wire Segment]
  \label{def:physics:phyx-mini-0966:phyxminiproblems-problemphyxmini0966-wiresegment}
  \lean{PhyXMiniProblems.ProblemPhyXMini0966.WireSegment}
  The four labelled pieces of the closed wire
\end{definition}
\begin{proof}
  This is the declared model definition of Wire Segment.
\end{proof}
\begin{definition}[Arc Segment]
  \label{def:physics:phyx-mini-0966:phyxminiproblems-problemphyxmini0966-arcsegment}
  \lean{PhyXMiniProblems.ProblemPhyXMini0966.ArcSegment}
  The two circular pieces, distinguished by radius and endpoint labels
\end{definition}
\begin{proof}
  This is the declared model definition of Arc Segment.
\end{proof}
\begin{definition}[Radial Connector]
  \label{def:physics:phyx-mini-0966:phyxminiproblems-problemphyxmini0966-radialconnector}
  \lean{PhyXMiniProblems.ProblemPhyXMini0966.RadialConnector}
  The two straight radial pieces
\end{definition}
\begin{proof}
  This is the declared model definition of Radial Connector.
\end{proof}
\begin{definition}[wire Segment Of Arc]
  \label{def:physics:phyx-mini-0966:phyxminiproblems-problemphyxmini0966-wiresegmentofarc}
  \lean{PhyXMiniProblems.ProblemPhyXMini0966.wireSegmentOfArc}
  \uses{def:physics:phyx-mini-0966:phyxminiproblems-problemphyxmini0966-wiresegment, def:physics:phyx-mini-0966:phyxminiproblems-problemphyxmini0966-arcsegment}
  Convert an arc label to the corresponding complete-wire segment label
\end{definition}
\begin{proof}
  This is the declared model definition of wire Segment Of Arc.
\end{proof}
\begin{definition}[wire Segment Of Connector]
  \label{def:physics:phyx-mini-0966:phyxminiproblems-problemphyxmini0966-wiresegmentofconnector}
  \lean{PhyXMiniProblems.ProblemPhyXMini0966.wireSegmentOfConnector}
  \uses{def:physics:phyx-mini-0966:phyxminiproblems-problemphyxmini0966-wiresegment, def:physics:phyx-mini-0966:phyxminiproblems-problemphyxmini0966-radialconnector}
  Convert a radial-connector label to its complete-wire segment label
\end{definition}
\begin{proof}
  This is the declared model definition of wire Segment Of Connector.
\end{proof}
\begin{definition}[Arc Current Sense]
  \label{def:physics:phyx-mini-0966:phyxminiproblems-problemphyxmini0966-arccurrentsense}
  \lean{PhyXMiniProblems.ProblemPhyXMini0966.ArcCurrentSense}
  Sense in which current traverses an arc as viewed in the page
\end{definition}
\begin{proof}
  This is the declared model definition of Arc Current Sense.
\end{proof}
\begin{definition}[Page Normal Direction]
  \label{def:physics:phyx-mini-0966:phyxminiproblems-problemphyxmini0966-pagenormaldirection}
  \lean{PhyXMiniProblems.ProblemPhyXMini0966.PageNormalDirection}
  The two directions normal to the plane of the image
\end{definition}
\begin{proof}
  This is the declared model definition of Page Normal Direction.
\end{proof}
\begin{definition}[Figure Color]
  \label{def:physics:phyx-mini-0966:phyxminiproblems-problemphyxmini0966-figurecolor}
  \lean{PhyXMiniProblems.ProblemPhyXMini0966.FigureColor}
  Color of the current arrows distinguished by the primary raster
\end{definition}
\begin{proof}
  This is the declared model definition of Figure Color.
\end{proof}
\begin{definition}[page Normal Vector]
  \label{def:physics:phyx-mini-0966:phyxminiproblems-problemphyxmini0966-pagenormalvector}
  \lean{PhyXMiniProblems.ProblemPhyXMini0966.pageNormalVector}
  \uses{def:physics:phyx-mini-0966:phyxminiproblems-problemphyxmini0966-spatialvector, def:physics:phyx-mini-0966:phyxminiproblems-problemphyxmini0966-pagenormaldirection}
  Unit normal associated with the two page-normal directions
\end{definition}
\begin{proof}
  This is the declared model definition of page Normal Vector.
\end{proof}
\begin{definition}[arc Field Normal Vector]
  \label{def:physics:phyx-mini-0966:phyxminiproblems-problemphyxmini0966-arcfieldnormalvector}
  \lean{PhyXMiniProblems.ProblemPhyXMini0966.arcFieldNormalVector}
  \uses{def:physics:phyx-mini-0966:phyxminiproblems-problemphyxmini0966-spatialvector, def:physics:phyx-mini-0966:phyxminiproblems-problemphyxmini0966-arccurrentsense, def:physics:phyx-mini-0966:phyxminiproblems-problemphyxmini0966-pagenormalvector}
  Right-hand-rule field direction at the center of a current-carrying arc
\end{definition}
\begin{proof}
  This is the declared model definition of arc Field Normal Vector.
\end{proof}
\begin{definition}[current Start Point]
  \label{def:physics:phyx-mini-0966:phyxminiproblems-problemphyxmini0966-currentstartpoint}
  \lean{PhyXMiniProblems.ProblemPhyXMini0966.currentStartPoint}
  \uses{def:physics:phyx-mini-0966:phyxminiproblems-problemphyxmini0966-figurepoint, def:physics:phyx-mini-0966:phyxminiproblems-problemphyxmini0966-wiresegment}
  Current-directed starting point of each labelled wire segment
\end{definition}
\begin{proof}
  This is the declared model definition of current Start Point.
\end{proof}
\begin{definition}[current End Point]
  \label{def:physics:phyx-mini-0966:phyxminiproblems-problemphyxmini0966-currentendpoint}
  \lean{PhyXMiniProblems.ProblemPhyXMini0966.currentEndPoint}
  \uses{def:physics:phyx-mini-0966:phyxminiproblems-problemphyxmini0966-figurepoint, def:physics:phyx-mini-0966:phyxminiproblems-problemphyxmini0966-wiresegment}
  Current-directed ending point of each labelled wire segment
\end{definition}
\begin{proof}
  This is the declared model definition of current End Point.
\end{proof}
\begin{definition}[Annular Sector Figure]
  \label{def:physics:phyx-mini-0966:phyxminiproblems-problemphyxmini0966-annularsectorfigure}
  \lean{PhyXMiniProblems.ProblemPhyXMini0966.AnnularSectorFigure}
  \uses{def:physics:phyx-mini-0966:phyxminiproblems-problemphyxmini0966-figurepoint, def:physics:phyx-mini-0966:phyxminiproblems-problemphyxmini0966-wiresegment, def:physics:phyx-mini-0966:phyxminiproblems-problemphyxmini0966-arcsegment, def:physics:phyx-mini-0966:phyxminiproblems-problemphyxmini0966-arccurrentsense, def:physics:phyx-mini-0966:phyxminiproblems-problemphyxmini0966-figurecolor}
  Literal point, segment, angle, and current-arrow evidence in image 966.png. The figure stores no magnetic-field value and no answer choice
\end{definition}
\begin{proof}
  This is the declared model definition of Annular Sector Figure.
\end{proof}
\begin{definition}[Annular Sector Wire Setup]
  \label{def:physics:phyx-mini-0966:phyxminiproblems-problemphyxmini0966-annularsectorwiresetup}
  \lean{PhyXMiniProblems.ProblemPhyXMini0966.AnnularSectorWireSetup}
  \uses{def:physics:phyx-mini-0966:phyxminiproblems-problemphyxmini0966-lengthquantity, def:physics:phyx-mini-0966:phyxminiproblems-problemphyxmini0966-electriccurrentquantity, def:physics:phyx-mini-0966:phyxminiproblems-problemphyxmini0966-magneticpermeabilityquantity, def:physics:phyx-mini-0966:phyxminiproblems-problemphyxmini0966-magneticfluxdensityquantity, def:physics:phyx-mini-0966:phyxminiproblems-problemphyxmini0966-figurepoint, def:physics:phyx-mini-0966:phyxminiproblems-problemphyxmini0966-wiresegment, def:physics:phyx-mini-0966:phyxminiproblems-problemphyxmini0966-arcsegment, def:physics:phyx-mini-0966:phyxminiproblems-problemphyxmini0966-radialconnector, def:physics:phyx-mini-0966:phyxminiproblems-problemphyxmini0966-annularsectorfigure}
  Independent physical quantities, geometry, and field observables. In particular, magneticFluxDensityMagnitudeAtP and totalMagneticField are not defined from the requested answer
\end{definition}
\begin{proof}
  This is the declared model definition of Annular Sector Wire Setup.
\end{proof}
\begin{definition}[segment Field Vector At PIn Teslas]
  \label{def:physics:phyx-mini-0966:phyxminiproblems-problemphyxmini0966-segmentfieldvectoratpinteslas}
  \lean{PhyXMiniProblems.ProblemPhyXMini0966.segmentFieldVectorAtPInTeslas}
  \uses{def:physics:phyx-mini-0966:phyxminiproblems-problemphyxmini0966-spatialvector, def:physics:phyx-mini-0966:phyxminiproblems-problemphyxmini0966-wiresegment, def:physics:phyx-mini-0966:phyxminiproblems-problemphyxmini0966-annularsectorwiresetup}
  SI vector field contribution of one segment, evaluated at P
\end{definition}
\begin{proof}
  This is the declared model definition of segment Field Vector At PIn Teslas.
\end{proof}
\begin{definition}[total Field Vector At PIn Teslas]
  \label{def:physics:phyx-mini-0966:phyxminiproblems-problemphyxmini0966-totalfieldvectoratpinteslas}
  \lean{PhyXMiniProblems.ProblemPhyXMini0966.totalFieldVectorAtPInTeslas}
  \uses{def:physics:phyx-mini-0966:phyxminiproblems-problemphyxmini0966-spatialvector, def:physics:phyx-mini-0966:phyxminiproblems-problemphyxmini0966-annularsectorwiresetup}
  Total SI magnetic-field vector evaluated at P
\end{definition}
\begin{proof}
  This is the declared model definition of total Field Vector At PIn Teslas.
\end{proof}
\begin{definition}[Matches Annular Sector Problem Data]
  \label{def:physics:phyx-mini-0966:phyxminiproblems-problemphyxmini0966-matchesannularsectorproblemdata}
  \lean{PhyXMiniProblems.ProblemPhyXMini0966.MatchesAnnularSectorProblemData}
  \uses{def:physics:phyx-mini-0966:phyxminiproblems-problemphyxmini0966-lengthincentimeters, def:physics:phyx-mini-0966:phyxminiproblems-problemphyxmini0966-currentinamperes, def:physics:phyx-mini-0966:phyxminiproblems-problemphyxmini0966-annularsectorwiresetup}
  Numerical quantities stated in the prose and displayed sector angle
\end{definition}
\begin{proof}
  This is the declared model definition of Matches Annular Sector Problem Data.
\end{proof}
\begin{definition}[Matches Primary Annular Sector Figure]
  \label{def:physics:phyx-mini-0966:phyxminiproblems-problemphyxmini0966-matchesprimaryannularsectorfigure}
  \lean{PhyXMiniProblems.ProblemPhyXMini0966.MatchesPrimaryAnnularSectorFigure}
  \uses{def:physics:phyx-mini-0966:phyxminiproblems-problemphyxmini0966-currentstartpoint, def:physics:phyx-mini-0966:phyxminiproblems-problemphyxmini0966-currentendpoint, def:physics:phyx-mini-0966:phyxminiproblems-problemphyxmini0966-annularsectorwiresetup}
  Primary-raster transcription. It fixes the closed traversal D → A → B → C → D, hence opposite arc senses, but asserts no field value
\end{definition}
\begin{proof}
  This is the declared model definition of Matches Primary Annular Sector Figure.
\end{proof}
\begin{definition}[Has Annular Sector Geometry]
  \label{def:physics:phyx-mini-0966:phyxminiproblems-problemphyxmini0966-hasannularsectorgeometry}
  \lean{PhyXMiniProblems.ProblemPhyXMini0966.HasAnnularSectorGeometry}
  \uses{def:physics:phyx-mini-0966:phyxminiproblems-problemphyxmini0966-lengthinmeters, def:physics:phyx-mini-0966:phyxminiproblems-problemphyxmini0966-annularsectorwiresetup}
  Metric and radial-ray content of the concentric annular-sector geometry. The scalar centralAngleRadians is the angle between the two radial rays
\end{definition}
\begin{proof}
  This is the declared model definition of Has Annular Sector Geometry.
\end{proof}
\begin{definition}[Has Physical Annular Sector Parameters]
  \label{def:physics:phyx-mini-0966:phyxminiproblems-problemphyxmini0966-hasphysicalannularsectorparameters}
  \lean{PhyXMiniProblems.ProblemPhyXMini0966.HasPhysicalAnnularSectorParameters}
  \uses{def:physics:phyx-mini-0966:phyxminiproblems-problemphyxmini0966-lengthinmeters, def:physics:phyx-mini-0966:phyxminiproblems-problemphyxmini0966-currentinamperes, def:physics:phyx-mini-0966:phyxminiproblems-problemphyxmini0966-permeabilityinteslametersperampere, def:physics:phyx-mini-0966:phyxminiproblems-problemphyxmini0966-annularsectorwiresetup}
  Positivity and branch assumptions for a nondegenerate annular sector
\end{definition}
\begin{proof}
  This is the declared model definition of Has Physical Annular Sector Parameters.
\end{proof}
\begin{definition}[Uses Standard Vacuum Permeability]
  \label{def:physics:phyx-mini-0966:phyxminiproblems-problemphyxmini0966-usesstandardvacuumpermeability}
  \lean{PhyXMiniProblems.ProblemPhyXMini0966.UsesStandardVacuumPermeability}
  \uses{def:physics:phyx-mini-0966:phyxminiproblems-problemphyxmini0966-permeabilityinteslametersperampere, def:physics:phyx-mini-0966:phyxminiproblems-problemphyxmini0966-annularsectorwiresetup}
  The physical permeability agrees with the selected Physlib electromagnetic system and has the standard vacuum SI value 4π × 10⁻⁷ T m/A
\end{definition}
\begin{proof}
  This is the declared model definition of Uses Standard Vacuum Permeability.
\end{proof}
\begin{definition}[Satisfies Steady Current Field Model]
  \label{def:physics:phyx-mini-0966:phyxminiproblems-problemphyxmini0966-satisfiessteadycurrentfieldmodel}
  \lean{PhyXMiniProblems.ProblemPhyXMini0966.SatisfiesSteadyCurrentFieldModel}
  \uses{def:physics:phyx-mini-0966:phyxminiproblems-problemphyxmini0966-annularsectorwiresetup}
  A steady line current produces time-independent segment and total fields
\end{definition}
\begin{proof}
  This is the declared model definition of Satisfies Steady Current Field Model.
\end{proof}
\begin{definition}[Satisfies Circular Arc Biot Savart Law]
  \label{def:physics:phyx-mini-0966:phyxminiproblems-problemphyxmini0966-satisfiescirculararcbiotsavartlaw}
  \lean{PhyXMiniProblems.ProblemPhyXMini0966.SatisfiesCircularArcBiotSavartLaw}
  \uses{def:physics:phyx-mini-0966:phyxminiproblems-problemphyxmini0966-lengthinmeters, def:physics:phyx-mini-0966:phyxminiproblems-problemphyxmini0966-currentinamperes, def:physics:phyx-mini-0966:phyxminiproblems-problemphyxmini0966-permeabilityinteslametersperampere, def:physics:phyx-mini-0966:phyxminiproblems-problemphyxmini0966-wiresegmentofarc, def:physics:phyx-mini-0966:phyxminiproblems-problemphyxmini0966-arcfieldnormalvector, def:physics:phyx-mini-0966:phyxminiproblems-problemphyxmini0966-annularsectorwiresetup, def:physics:phyx-mini-0966:phyxminiproblems-problemphyxmini0966-segmentfieldvectoratpinteslas}
  At their common center, each circular arc obeys B = μ₀ I θ / (4πR) with direction selected by the right-hand rule. This is a general governing law in the stored physical quantities, not the requested numerical result
\end{definition}
\begin{proof}
  This is the declared model definition of Satisfies Circular Arc Biot Savart Law.
\end{proof}
\begin{definition}[Satisfies Radial Connector Biot Savart Law]
  \label{def:physics:phyx-mini-0966:phyxminiproblems-problemphyxmini0966-satisfiesradialconnectorbiotsavartlaw}
  \lean{PhyXMiniProblems.ProblemPhyXMini0966.SatisfiesRadialConnectorBiotSavartLaw}
  \uses{def:physics:phyx-mini-0966:phyxminiproblems-problemphyxmini0966-wiresegmentofconnector, def:physics:phyx-mini-0966:phyxminiproblems-problemphyxmini0966-annularsectorwiresetup, def:physics:phyx-mini-0966:phyxminiproblems-problemphyxmini0966-segmentfieldvectoratpinteslas}
  A straight radial current element has dℓ ∥ r, so its Biot--Savart cross product and hence its field at the center P vanish
\end{definition}
\begin{proof}
  This is the declared model definition of Satisfies Radial Connector Biot Savart Law.
\end{proof}
\begin{definition}[Satisfies Magnetic Field Superposition]
  \label{def:physics:phyx-mini-0966:phyxminiproblems-problemphyxmini0966-satisfiesmagneticfieldsuperposition}
  \lean{PhyXMiniProblems.ProblemPhyXMini0966.SatisfiesMagneticFieldSuperposition}
  \uses{def:physics:phyx-mini-0966:phyxminiproblems-problemphyxmini0966-annularsectorwiresetup, def:physics:phyx-mini-0966:phyxminiproblems-problemphyxmini0966-segmentfieldvectoratpinteslas, def:physics:phyx-mini-0966:phyxminiproblems-problemphyxmini0966-totalfieldvectoratpinteslas}
  The total field at P is the vector sum of all four segment fields
\end{definition}
\begin{proof}
  This is the declared model definition of Satisfies Magnetic Field Superposition.
\end{proof}
\begin{definition}[Calibrates Magnetic Flux Density Magnitude At P]
  \label{def:physics:phyx-mini-0966:phyxminiproblems-problemphyxmini0966-calibratesmagneticfluxdensitymagnitudeatp}
  \lean{PhyXMiniProblems.ProblemPhyXMini0966.CalibratesMagneticFluxDensityMagnitudeAtP}
  \uses{def:physics:phyx-mini-0966:phyxminiproblems-problemphyxmini0966-magneticfluxdensityinteslas, def:physics:phyx-mini-0966:phyxminiproblems-problemphyxmini0966-annularsectorwiresetup, def:physics:phyx-mini-0966:phyxminiproblems-problemphyxmini0966-totalfieldvectoratpinteslas}
  The independent dimensionful magnitude observable is calibrated by the norm of the total Physlib magnetic-field vector at P
\end{definition}
\begin{proof}
  This is the declared model definition of Calibrates Magnetic Flux Density Magnitude At P.
\end{proof}
\begin{definition}[Answer Choice]
  \label{def:physics:phyx-mini-0966:phyxminiproblems-problemphyxmini0966-answerchoice}
  \lean{PhyXMiniProblems.ProblemPhyXMini0966.AnswerChoice}
  Labels of the four displayed magnetic-field magnitudes
\end{definition}
\begin{proof}
  This is the declared model definition of Answer Choice.
\end{proof}
\begin{definition}[displayed Magnetic Field In Microteslas]
  \label{def:physics:phyx-mini-0966:phyxminiproblems-problemphyxmini0966-displayedmagneticfieldinmicroteslas}
  \lean{PhyXMiniProblems.ProblemPhyXMini0966.displayedMagneticFieldInMicroteslas}
  \uses{def:physics:phyx-mini-0966:phyxminiproblems-problemphyxmini0966-answerchoice}
  Magnetic-field magnitude printed beside each choice, in microteslas
\end{definition}
\begin{proof}
  This is the declared model definition of displayed Magnetic Field In Microteslas.
\end{proof}
\begin{definition}[Rounds To Nearest Hundredth Microtesla]
  \label{def:physics:phyx-mini-0966:phyxminiproblems-problemphyxmini0966-roundstonearesthundredthmicrotesla}
  \lean{PhyXMiniProblems.ProblemPhyXMini0966.RoundsToNearestHundredthMicrotesla}
  Rounding to the nearest 0.01 μT, the precision of every answer choice
\end{definition}
\begin{proof}
  This is the declared model definition of Rounds To Nearest Hundredth Microtesla.
\end{proof}
\begin{definition}[Is Unique Closest Displayed Answer]
  \label{def:physics:phyx-mini-0966:phyxminiproblems-problemphyxmini0966-isuniqueclosestdisplayedanswer}
  \lean{PhyXMiniProblems.ProblemPhyXMini0966.IsUniqueClosestDisplayedAnswer}
  \uses{def:physics:phyx-mini-0966:phyxminiproblems-problemphyxmini0966-answerchoice, def:physics:phyx-mini-0966:phyxminiproblems-problemphyxmini0966-displayedmagneticfieldinmicroteslas}
  A displayed answer is strictly closer than every distinct alternative
\end{definition}
\begin{proof}
  This is the declared model definition of Is Unique Closest Displayed Answer.
\end{proof}
\begin{lemma}[total Magnetic Field At P exact]
  \label{lem:physics:phyx-mini-0966:phyxminiproblems-problemphyxmini0966-totalmagneticfieldatp-exact}
  \lean{PhyXMiniProblems.ProblemPhyXMini0966.totalMagneticFieldAtP_exact}
  \uses{def:physics:phyx-mini-0966:phyxminiproblems-problemphyxmini0966-pagenormalvector, def:physics:phyx-mini-0966:phyxminiproblems-problemphyxmini0966-annularsectorwiresetup, def:physics:phyx-mini-0966:phyxminiproblems-problemphyxmini0966-totalfieldvectoratpinteslas, def:physics:phyx-mini-0966:phyxminiproblems-problemphyxmini0966-matchesannularsectorproblemdata, def:physics:phyx-mini-0966:phyxminiproblems-problemphyxmini0966-matchesprimaryannularsectorfigure, def:physics:phyx-mini-0966:phyxminiproblems-problemphyxmini0966-hasannularsectorgeometry, def:physics:phyx-mini-0966:phyxminiproblems-problemphyxmini0966-hasphysicalannularsectorparameters, def:physics:phyx-mini-0966:phyxminiproblems-problemphyxmini0966-usesstandardvacuumpermeability, def:physics:phyx-mini-0966:phyxminiproblems-problemphyxmini0966-satisfiessteadycurrentfieldmodel, def:physics:phyx-mini-0966:phyxminiproblems-problemphyxmini0966-satisfiescirculararcbiotsavartlaw, def:physics:phyx-mini-0966:phyxminiproblems-problemphyxmini0966-satisfiesradialconnectorbiotsavartlaw, def:physics:phyx-mini-0966:phyxminiproblems-problemphyxmini0966-satisfiesmagneticfieldsuperposition}
  The inner clockwise arc dominates the outer counterclockwise arc. Their vector difference is exactly (4π/3) μT into the page; both radial connector contributions vanish
\end{lemma}
\begin{proof}
  The conclusion follows from the governing relations and figure readouts named in the statement of total Magnetic Field At P exact.
\end{proof}
% --- Archon declaration topology end ---
% --- Archon physics formalization source end ---
